\documentclass[11pt]{article}
\usepackage[a4paper,margin=25mm]{geometry}
\usepackage{amsmath,amssymb,amsthm,booktabs,microtype}
\usepackage[T1]{fontenc}
\usepackage{lmodern}
\usepackage[dvipsnames]{xcolor}
\usepackage[colorlinks=true,linkcolor=MidnightBlue,citecolor=MidnightBlue,urlcolor=MidnightBlue]{hyperref}
\usepackage{fancyhdr}
\pagestyle{fancy}\fancyhf{}
\fancyhead[L]{\small Order-polynomial positivity with a short chain}
\fancyhead[R]{\small Dev D. Goyal}
\fancyfoot[C]{\thepage}
\setlength{\headheight}{14pt}
\hypersetup{pdftitle={Order-polynomial positivity for width-two posets with a chain of at most five elements},
pdfauthor={Dev D. Goyal},
pdfsubject={Fully autonomous agentic research},
pdfkeywords={order polynomial, width-two posets, Ehrhart positivity, order polytope, non-nesting rooks, sum of squares, Lean 4}}
\newtheorem{theorem}{Theorem}
\newtheorem{lemma}[theorem]{Lemma}
\newtheorem{proposition}[theorem]{Proposition}
\newtheorem{corollary}[theorem]{Corollary}
\theoremstyle{definition}
\newtheorem{definition}[theorem]{Definition}
\theoremstyle{remark}
\newtheorem{remark}[theorem]{Remark}
\newtheorem{example}[theorem]{Example}
\newcommand{\rise}[2]{#1^{\overline{#2}}}
\title{\textbf{Order-polynomial positivity for width-two posets\\ with a chain of at most five elements}}
\author{Dev D. Goyal}
\date{31 October 2026}
\begin{document}
\maketitle
\begin{abstract}
For a finite poset $P$ let $\Omega_P(t)$ be its order polynomial, which
counts weakly order-preserving maps $P\to\{1,\dots,m\}$ at $t=m$.
Ferroni, Morales and Panova conjectured that $|P|!\,\Omega_P(t)$ has
nonnegative coefficients in the monomial basis whenever $P$ has width at
most two, and reported computer checks when both chains of a two-chain
cover have at most five elements. We prove the conjecture for every
finite poset that is the union of two chains $A$ and $B$ with $|A|\le5$,
the chain $B$ being arbitrarily long and the comparabilities between $A$
and $B$ arbitrary. For $|A|,|B|\ge1$ we also prove
$[t]\,\Omega_P(t)\ge(|A|-1)!\,(|B|-1)!\,m/|P|!$, where $m$ is the number of
elements of $B$ comparable with some element of $A$, so that the linear
coefficient is positive exactly when the comparability graph is
connected. Via the descent--rook correspondence of Alexandersson and
Jal, the normalized linear coefficient is, on each of the $64$ chambers
of endpoint configurations, a polynomial of degree at most $|A|$ in the
gaps between endpoints, independent of $|B|$. Exact certificates (sums
of $1{,}072$ rational squares plus nonnegative remainders, matched with
an independent order-polynomial evaluation at $194{,}112$ determining
points) prove these polynomials nonnegative, and the known
ideal--filter reduction of Ferroni, Morales and Panova lifts the linear
coefficient to all coefficients. Lean~4 checks the nonnegativity
certificates of the chamber polynomials for $|A|\le4$; the rest is a
conventional proof with certificates checked by an independent program.
\end{abstract}

\section{Introduction and results}

Let $P$ be a finite poset with $n$ elements. For a positive integer $m$
let $\Omega_P(m)$ be the number of maps $f:P\to\{1,\dots,m\}$ with
$f(x)\le f(y)$ whenever $x<y$. It is a polynomial in $m$ of degree $n$,
the \emph{order polynomial}, and $\Omega_P(t+1)$ is the Ehrhart
polynomial of the order polytope of $P$~\cite{Stanley}. Writing
$\Omega_P(t)=\sum_k e_k\binom tk$, where $e_k$ counts chains
$\emptyset=I_0\subsetneq I_1\subsetneq\dots\subsetneq I_k=P$ of order
ideals, shows that $n!\,\Omega_P(t)\in\mathbb Z[t]$. Nonnegativity of
the coefficients of $\Omega_P(t)$ implies that of $\Omega_P(t+1)$, hence
Ehrhart positivity of the order polytope, but it fails in general: the
star with one minimum and four maximal elements has
$\Omega(t)=\tfrac{t^5}5+\tfrac{t^4}2+\tfrac{t^3}3-\tfrac t{30}$.

\paragraph{The conjecture.} Ferroni, Morales and Panova~\cite[Section~7.1,
Conjecture~7.1]{FMP} conjectured that $n!\,\Omega_P(t)$ has nonnegative
coefficients in the monomial basis for every poset $P$ of width at most
two. By Dilworth's theorem these are exactly the posets covered by two
chains. They reported computer verification when both chains have at
most five elements, proved positivity for the cell posets of skew shapes
(a different construction of a poset from a shape), and showed that
for a family closed under order ideals and their complementary filters
positivity of all coefficients follows from positivity of the linear
coefficient~\cite[Proposition~3.3, Theorem~3.4]{FMP}.

\begin{theorem}\label{thm:main}
Let $P$ be a finite poset that is the disjoint union of two chains $A$ and
$B$ with $|A|\le5$. There is no condition on $|B|$ or on the
comparabilities between $A$ and $B$. Then every coefficient of
$|P|!\,\Omega_P(t)$ in the monomial basis is a nonnegative integer.
\end{theorem}

\begin{theorem}\label{thm:lb}
In Theorem~\ref{thm:main} let $a=|A|\ge1$, $b=|B|\ge1$ and let $m_A(B)$
be the number of elements of $B$ comparable with at least one element
of $A$. Then
\[
[t]\,\Omega_P(t)\ \ge\ \frac{(a-1)!\,(b-1)!}{(a+b)!}\,m_A(B).
\]
\end{theorem}

\begin{corollary}\label{cor:conn}
For a nonempty poset $P$ as in Theorem~\ref{thm:main}, the linear
coefficient of $\Omega_P(t)$ is positive if the comparability graph of
$P$ is connected and zero otherwise.
\end{corollary}

The empty poset ($\Omega=1$) and the case of an empty chain are included.
The hypothesis is a partition into two chains with one of them short;
the existence of a long chain whose complement has at most five elements
is weaker, since the complement need not be a chain. The case $a=1$ is
elementary (Example~\ref{ex:a1}). Five is the limit of the certificates
we have, not a proven threshold.

\paragraph{Method.} Section~\ref{sec:enc} encodes the comparabilities
between the chains by two monotone boundary sequences. Using the
correspondence between descents of linear extensions and non-nesting
rook placements due to Alexandersson and Jal~\cite{AJ}, Section~\ref{sec:lin}
writes $n!\,[t]\Omega_P(t)/(b-1)!$ as an alternating combination of rook
numbers. Section~\ref{sec:chambers} splits the boundary configurations
into $C_a$ chambers, $C_a$ the Catalan number, and shows that on each
chamber this functional is a polynomial of degree at most $a$ in the
$2a+1$ gaps between consecutive endpoints, with no dependence on $b$
beyond the gaps. Finitely many exact evaluations therefore identify it,
and Section~\ref{sec:cert} certifies each chamber polynomial nonnegative
on the orthant. For $a\le3$ all its coefficients are positive; for
$a=4$, ten of the fourteen chamber polynomials have negative
coefficients and are written as sums of rational squares plus a
remainder with nonnegative coefficients. Section~\ref{sec:proof} applies
the ideal--filter reduction.

\paragraph{Related work.} Section~\ref{sec:related} compares the result
with two public partial results on the same conjecture: positivity for
ordinal sums of one- and two-element antichains, and a coefficient
positivity theorem for lattice-point counts whose order-polynomial
factors have a single cross-link. Neither class contains the class of
Theorem~\ref{thm:main}, nor is contained in it.

\section{Boundary encoding}\label{sec:enc}

Write $A_1<\dots<A_a$ and $B_1<\dots<B_b$. For $1\le i\le a$ let $L_i$ be
the number of elements of $B$ below $A_i$ and $U_i$ be $b$ minus the
number of elements of $B$ above $A_i$. By transitivity within the chains,
\begin{equation}\label{eq:mono}
0\le L_1\le\dots\le L_a\le b,\qquad 0\le U_1\le\dots\le U_a\le b,\qquad
L_i\le U_i,
\end{equation}
and the cross relation is
\begin{equation}\label{eq:cross}
B_j<A_i\iff j\le L_i,\qquad A_i<B_j\iff j>U_i,\qquad A_i\parallel
B_j\iff L_i<j\le U_i.
\end{equation}
Conversely, any data satisfying~\eqref{eq:mono}, with~\eqref{eq:cross}
and the two chain orders, define a partial order. Antisymmetry holds since
$j\le L_i$ and $j>U_i$ are incompatible. For transitivity through the
other chain: $A_i<B_j<A_{i'}$ gives $U_i<j\le L_{i'}\le U_{i'}$, hence
$i<i'$ by monotonicity of $U$; $B_j<A_i<B_{j'}$ gives $j\le L_i\le U_i<j'$;
the remaining patterns combine one cross relation with a chain relation
and follow from monotonicity. Every pair $(L,U)$ satisfying~\eqref{eq:mono}
therefore arises, and the incomparable pairs form the board of cells
$(i,j)$ with $L_i<j\le U_i$.

\section{The linear coefficient as a rook functional}\label{sec:lin}

A \emph{non-nesting rook placement} of size $k$ is a set of $k$
incomparable pairs $(A_{i_1},B_{j_1}),\dots,(A_{i_k},B_{j_k})$ with
$i_1<\dots<i_k$ and $j_1<\dots<j_k$; let $r_k$ be their number, $r_0=1$.

\begin{lemma}[Alexandersson--Jal]\label{lem:rook}
Label $P$ naturally by a linear extension that, whenever both are
available, takes an element of $A$ before an element of $B$. Then the
number of linear extensions of $P$ with $k$ descents is $r_k$.
\end{lemma}

\begin{proof}
This is the correspondence of~\cite[Theorem~5.4, Lemma~5.6,
Corollary~5.8]{AJ}; we give the direct argument covering empty rows,
columns and coincident endpoints. If $A_i\parallel B_j$ then
$L_i<j$, and by induction on $i$ every element below $A_i$ precedes $B_j$
in the labelling, so
$A_i$ is available before $B_j$ is chosen and receives the smaller
label. A descent of a linear extension is a pair of adjacent incomparable
elements in decreasing label order, hence a pair $B_jA_i$ with
$A_i\parallel B_j$, and the recorded pairs of an extension have
increasing indices in both chains. Conversely, given a non-nesting
placement, the prefixes $\{A_1,\dots,A_{i_s-1}\}\cup\{B_1,\dots,B_{j_s-1}\}$
and $\{A_1,\dots,A_{i_s}\}\cup\{B_1,\dots,B_{j_s}\}$ are order ideals by
$L_{i_s}<j_s\le U_{i_s}$ and~\eqref{eq:mono}; filling the elements between
consecutive prescribed pairs in increasing label order gives the unique
linear extension whose descents are exactly the prescribed pairs.
\end{proof}

By the standard descent expansion~\cite[Section~3.15]{Stanley}, for a
natural labelling,
\begin{equation}\label{eq:desc}
\Omega_P(t)=\sum_{k=0}^{a}r_k\binom{t+n-1-k}{n},\qquad n=a+b.
\end{equation}
Assume $a,b\ge1$. For $0\le k\le n-1$ the binomial
$\binom{t+n-1-k}{n}$ is $(t+n-1-k)\cdots(t-k)/n!$, whose linear
coefficient is $(-1)^kk!\,(n-1-k)!/n!$. Since $k\le a\le n-1$,
\begin{equation}\label{eq:F}
F:=\frac{n!\,[t]\,\Omega_P(t)}{(b-1)!}=\sum_{k=0}^{a}(-1)^k\,k!\,\rise b{a-k}\,r_k,
\qquad \rise bd=b(b+1)\cdots(b+d-1).
\end{equation}
This is the linear coefficient of the unshifted order polynomial, not of
the Ehrhart polynomial $\Omega_P(t+1)$.

\section{Chambers and gap polynomials}\label{sec:chambers}

Merge the lists $L_1,\dots,L_a$ and $U_1,\dots,U_a$ into one weakly
increasing list, placing $L$-endpoints before $U$-endpoints at ties and
keeping the order within each list. Recording the endpoint types gives a
word $w$ with $a$ letters $L$ and $a$ letters $U$; since $L_i\le U_i$,
every prefix has at least as many $L$'s as $U$'s, so $w$ is a Dyck word.
There are $C_a=\frac1{a+1}\binom{2a}a$ of them: $1,2,5,14,42$ for
$a=1,\dots,5$, in total $64$ chambers. Let $x_0$ be the first endpoint,
$x_j$ the distance from the $j$-th to the $(j+1)$-st endpoint for
$1\le j<2a$, and $x_{2a}=b$ minus the last endpoint. Then
$x\in\mathbb Z_{\ge0}^{2a+1}$, $b=x_0+\dots+x_{2a}$, and conversely every
pair $(w,x)$ gives boundary data satisfying~\eqref{eq:mono}, hence a
poset. Ties are zero gaps; a configuration on a chamber boundary may
have several encodings, which is harmless.

If after $j$ endpoints $o$ letters $L$ and $c$ letters $U$ have
occurred, every element $B_{j'}$ in the $j$-th gap segment is
incomparable exactly with $A_{c+1},\dots,A_o$, by~\eqref{eq:cross}. Write
$h_j=o-c$ for the number of these rows; $h_0=h_{2a}=0$.

\begin{lemma}\label{lem:poly}
For fixed $(a,w)$ there is a polynomial $F_{a,w}\in\mathbb Q[x_0,\dots,x_{2a}]$
of total degree at most $a$, with $F_{a,w}(0)=0$, such that $F$
in~\eqref{eq:F} equals $F_{a,w}(x)$ for every poset encoded by $(w,x)$
with $b\ge1$.
\end{lemma}

\begin{proof}
A placement of size $k$ selects $k$ rows and assigns them, in increasing
order, to weakly increasing gap segments in which they are active; if
segment $j$ receives $d_j$ rows, their columns can be chosen in
$\binom{x_j}{d_j}$ ways, and the columns are then increasing. Hence
$r_k=\sum\prod_j\binom{x_j}{d_j}$ over the finitely many such
assignments, a polynomial of degree at most $k$. Since $b=\sum x_j$,
each term of~\eqref{eq:F} has degree at most $a$. At $x=0$ every $r_k$
with $k\ge1$ vanishes and $\rise 0a=0$.
\end{proof}

\begin{lemma}\label{lem:det}
A polynomial of total degree at most $a$ in $d$ variables that vanishes on
$S_{d,a}=\{x\in\mathbb Z_{\ge0}^d:x_0+\dots+x_{d-1}\le a\}$ is zero.
\end{lemma}

\begin{proof}
The products $\prod_i\binom{x_i}{\alpha_i}$ with $|\alpha|\le a$ form a
basis. Evaluating at $x=\beta\in S_{d,a}$ in order of increasing
$|\beta|$, only $\alpha\le\beta$ contribute, the coefficients with
$|\alpha|<|\beta|$ are already zero, and among $|\alpha|=|\beta|$ only
$\alpha=\beta$ survives, with value one.
\end{proof}

Thus $F_{a,w}$ is determined by its values on $S_{2a+1,a}$, which has
$\binom{3a+1}{a}$ points: $4,21,120,715,4368$ for $a=1,\dots,5$.

\section{Certificates}\label{sec:cert}

For each of the $64$ chambers the certificate supplies $F_{a,w}$
explicitly by its monomial coefficients, and a proof that it is
nonnegative on $\mathbb R_{\ge0}^{2a+1}$:
\begin{itemize}
\item for $a\le3$, all coefficients of $F_{a,w}$ are positive integers;
\item for $a=4,5$, an identity
\begin{equation}\label{eq:sos}
F_{a,w}(x)=\sum_s q_s\,x^{\mu_s}\bigl(\rho_s x^{u_s}-x^{v_s}\bigr)^2+\sum_e c_e\,x^e
\end{equation}
with rational $q_s\ge0$, $c_e\ge0$, rational $\rho_s$ and exponent vectors
$\mu_s,u_s,v_s,e\in\mathbb Z_{\ge0}^{2a+1}$; every term is nonnegative on
the orthant.
\end{itemize}

\begin{table}[h]
\centering
\begin{tabular}{rrrrr}
\toprule
$a$ & chambers & points per chamber & determining points & square terms\\
\midrule
1 & 1 & 4 & 4 & 0\\
2 & 2 & 21 & 42 & 0\\
3 & 5 & 120 & 600 & 0\\
4 & 14 & 715 & 10{,}010 & 101\\
5 & 42 & 4{,}368 & 183{,}456 & 971\\
\midrule
total & 64 & & 194{,}112 & 1{,}072\\
\bottomrule
\end{tabular}
\caption{Size of the certificate.}\label{tab:cert}
\end{table}

The checker \texttt{certificate/check\_certificates.py}, which uses only
the Python standard library and exact rational arithmetic, verifies:
\begin{enumerate}
\item the chamber list for each $a\le5$ is exactly the set of Dyck words;
\item each $F_{a,w}$ has $2a+1$ variables, total degree at most $a$ and
zero constant term;
\item positivity of the coefficients for $a\le3$; for $a=4,5$
nonnegativity of every $q_s$ and $c_e$ and exact coefficientwise
equality in~\eqref{eq:sos};
\item at every point $x\in S_{2a+1,a}$, equality of $F_{a,w}(x)$ with the
value of~\eqref{eq:F} for the poset encoded by $(w,x)$, computed from
scratch: $\Omega_P(m)$ for $m=0,\dots,n$ by counting multichains of
order ideals, which are pairs of chain prefixes, and the linear
coefficient from finite differences; this evaluation does not use rook
numbers;
\item the degree-one remainder coefficients needed in
Section~\ref{sec:lb}.
\end{enumerate}
As a sanity check of Lemma~\ref{lem:poly}, the checker also compares
$F_{a,w}$ with~\eqref{eq:F} at three random gap vectors with $|x|>a$ in
every chamber, and it rejects corrupted certificates, among them a
polynomial and remainder shifted by the same monomial, which still
satisfy~\eqref{eq:sos} and are caught only by the point identities. The certificates and the identities were found
with the help of a linear-programming search; the search is not part of
the proof, and the accepted weights are exact rationals. The original
research package contained two further independent verifications of the
same identities, which we reran.

By Lemmas~\ref{lem:poly} and~\ref{lem:det} the identities hold as
polynomial identities, so $F_{a,w}(x)\ge0$ for every $x\ge0$, and
by~\eqref{eq:F} $[t]\,\Omega_P(t)\ge0$ for every $P$ in
Theorem~\ref{thm:main} with $a,b\ge1$.

\begin{remark}
For $a=4$ the coefficientwise argument is not available: ten of the
fourteen polynomials $F_{4,w}$ have a negative coefficient. A negative
coefficient of a gap polynomial is not a negative value of an order
polynomial. For $a=6$ two limited searches for certificates of the
form~\eqref{eq:sos} failed; this says nothing about the truth of the
conjecture there.
\end{remark}

\section{Proof of Theorem~\ref{thm:main}}\label{sec:proof}

The class $\mathcal S_5$ of posets admitting a partition into two chains,
one of size at most five, is closed under order ideals and their
complementary filters: intersect the two chains with the ideal or
filter. For a nonempty poset $P$ write $\Omega_P(t)=\sum_{k\ge1}c_k(P)t^k$;
$\Omega_P(0)=0$. Splitting a map to $\{1,\dots,x+y\}$ at the largest
value at most $x$ gives
$\Omega_P(x+y)=\sum_I\Omega_I(x)\,\Omega_{P\setminus I}(y)$ over order
ideals $I$; putting $x=y=t$ and comparing coefficients,
\begin{equation}\label{eq:red}
(2^k-2)\,c_k(P)=\sum_{\emptyset\ne I\ne P}\ \sum_{i=1}^{k-1}c_i(I)\,c_{k-i}(P\setminus I)\qquad(k\ge2).
\end{equation}
This is the argument of~\cite[Proposition~3.3, Theorem~3.4]{FMP}. By
induction on $|P|$ the right-hand side is nonnegative for $P\in\mathcal
S_5$, and $c_1(P)\ge0$ by Section~\ref{sec:cert} when both chains are
nonempty. If a chain is empty, $P$ is a chain and
$\Omega_P(t)=t(t+1)\cdots(t+n-1)/n!$. Integrality was noted in the
introduction. \qed

\section{The linear lower bound}\label{sec:lb}

The checker also verifies that in every chamber the degree-one part of
the nonnegative part of the certificate (the polynomial itself for
$a\le3$, the remainder in~\eqref{eq:sos} for $a=4,5$) is exactly
\[
(a-1)!\sum_{j:\,h_j<a}x_j .
\]
All other terms are nonnegative on the orthant, so
$F_{a,w}(x)\ge(a-1)!\sum_{h_j<a}x_j$. An element of $B$ in gap segment
$j$ is incomparable with exactly $h_j$ elements of $A$, so it is
comparable with some element of $A$ if and only if $h_j<a$, and the sum
equals $m_A(B)$. Multiplying by $(b-1)!/n!$ proves Theorem~\ref{thm:lb}.

For Corollary~\ref{cor:conn}: if $a=0$ or $b=0$ the poset is a nonempty
chain, which is connected and has linear coefficient $1/n$. Otherwise,
if some element of $B$ is comparable with some element of $A$ then
$m_A(B)\ge1$, the comparability graph is connected because each chain
is, and the linear coefficient is positive. If there is no such
comparability, $P$ is the disjoint union of two nonempty chains,
$\Omega_P$ is the product of their order polynomials, each without
constant term, and the linear coefficient is zero. \qed

\begin{example}\label{ex:a1}
For $a=1$ with boundaries $L\le U$ there are $U-L$ incomparable elements,
$r_1=U-L$ and $F=b-(U-L)=m_A(B)$, so Theorem~\ref{thm:lb} is an equality:
$[t]\,\Omega_P(t)=m_A(B)/(b(b+1))$. For $b=8$, $L=2$, $U=6$ it is $1/18$.
\end{example}

\begin{example}
For $a=5$, $b=20$, $L=(0,2,5,5,11)$ and $U=(4,8,12,17,20)$, every element
of $B$ is comparable with some element of $A$, and Theorem~\ref{thm:lb}
gives $[t]\,\Omega_P(t)\ge4!\,19!\,20/25!=1/265650$; the exact value is
$7123/531300$. This poset is outside the range of the computations
in~\cite{FMP}. For $L=(0,\dots,0)$ and $U=(20,\dots,20)$ there are no
cross-comparabilities and $\Omega_P(t)=\binom{t+4}5\binom{t+19}{20}$ has
zero linear coefficient.
\end{example}

\section{Formal verification}\label{sec:lean}

The Lean~4 project (Lean~4.35.0, Mathlib
\texttt{56a005855ea17e0b0958f73ec801ae140b0a1679}) checks the chamber
certificates for $a\le4$. A chamber is a record of the word $w$, a
scale $D$ and integer data: the polynomial $D\,F_{a,w}$, the degree-one
part and the rest of $D$ times the nonnegative remainder, and square
triples $(p,q,r)$ standing for $p\,x^{\mu+2u}-q\,x^{\mu+u+v}+r\,x^{\mu+2v}$.
The file \texttt{ShortChain/Data.lean} is generated from the certificate
files by \texttt{scripts/generate\_lean\_data.py}, which also checks
that it matches them. The Boolean test \texttt{Chamber.valid} is
evaluated by the kernel (\texttt{decide +kernel}) for all $22$ chambers:
the remainder and the expanded squares add up to the polynomial, the
remainder coefficients and $p,r$ are nonnegative and $q^2\le4pr$, and
the degree-one part is $D\,(a-1)!$ exactly on the segments with
$h_j<a$. The theorem \texttt{ShortChain.certified\_chambers} states that
for every chamber in the list and every gap vector $x$ with nonnegative
entries in any linearly ordered field,
\[
0\le F_{a,w}(x)\quad\text{and}\quad (a-1)!\sum_{h_j<a}x_j\le F_{a,w}(x),
\]
and \texttt{ShortChain.chambers\_complete} that for $a=1,\dots,4$ the
chamber words are exactly the Dyck words. No \texttt{native\_decide}
is used, and a module-level audit shows that every declaration depends
only on \texttt{propext}, \texttt{Classical.choice} and
\texttt{Quot.sound}.

Not formalized are the identification of $F_{a,w}$ with the linear
coefficient of the order polynomial (Lemmas~\ref{lem:rook}
and~\ref{lem:poly}, the descent expansion~\eqref{eq:desc}, the
determining-set Lemma~\ref{lem:det} and the point identities), the
chambers with $a=5$, whose $177{,}176$ polynomial terms are checked only
by the Python program, and the reduction of Section~\ref{sec:proof}.

\section{Related work and scope}\label{sec:related}

The conjecture, the reduction to the linear coefficient and the
computations for both chains of size at most five are
from~\cite{FMP}; the descent--rook correspondence is
from~\cite{AJ}, and the descent expansion is classical~\cite{Stanley}.
Liu, Xin and Zhang~\cite{LXZ} proved Ehrhart positivity of order
polytopes of dimension at most $13$, which concerns $\Omega_P(t+1)$ and
bounded $|P|$. Jochemko and Menon~\cite{JM} showed that for an
ideal--filter closed family with nonnegative linear coefficients, marked
order polynomials have nonnegative coefficients in the marking gaps;
Theorem~\ref{thm:main} supplies that hypothesis for $\mathcal S_5$.
Kahane~\cite{Kahane} interprets the coefficients for fences.

Two public partial results on Conjecture~7.1 of~\cite{FMP} were found
while preparing this note. The repository
\href{https://github.com/10m-u/open-problems-catalog/tree/ca8d4d44ba3db527f19f28928c240168f545ecfe}{10m-u/open-problems-catalog}
(commit of 14 October 2026) proves nonnegativity for all ordinal sums of
antichains of size one or two, by a root-location argument. Such a poset
lies in $\mathcal S_5$ exactly when it has at most five two-element
layers, while $\mathcal S_5$ contains posets outside the closure of
these layers under ordinal and disjoint sums, the smallest being the
four-element poset with relations $a<b$, $c<b$, $c<d$, whose
comparability and incomparability graphs are both connected. The
repository
\href{https://github.com/chris-hadad/lr-stretching-results/blob/0677a8b4bda0db770ad690824c236222783eeb95/results/overlap-families/proofs/A58-COUPLED-PROOF.md}{chris-hadad/lr-stretching-results}
(commit of 29 October 2026) proves coefficient positivity for a
lattice-point count whose order-polynomial factors are two-chain posets
with a single cross-link; those factors may have both chains longer than
five, but their cross-relations are far more restricted than in
$\mathcal S_5$. Neither result implies Theorem~\ref{thm:main}, and
Theorem~\ref{thm:main} implies neither. The present theorem is not the
first all-size partial result on the conjecture. Searches in October
2026 found no statement of Theorem~\ref{thm:main} or~\ref{thm:lb};
coverage of recent and unpublished work is incomplete and no priority
is claimed.

The unrestricted conjecture remains open. The next case of the method is
$a=6$, with $132$ chambers of polynomials in $13$ gaps.

\section*{Research process and AI disclosure}

This is fully autonomous agentic research. Dev D. Goyal is responsible
for publication. No expert validation or external peer review is
claimed, and the novelty assessment is provisional.

The theorem, the chamber reduction, the certificates and their
verifiers were produced by autonomous AI research agents in a GPT-6
Astra research workflow, according to its records. The workflow's
internal reviews, including a separate referee reconstruction of the
rook polynomials, were AI-assisted reviews of the proof and not peer
review. Its own late audit found the two public partial results
discussed in Section~\ref{sec:related}.

Opus 5.5 by Anthropic used with Claude Code assessed the package for this
release. It checked the proof, wrote an independent certificate checker
with its own order-polynomial evaluation, reran the original
verifications, compared the statement with the cited sources and the two
public partial results, wrote and ran the Lean formalization described in
Section~\ref{sec:lean}, and wrote this note. AI checks are not human peer
review.

\begin{thebibliography}{9}
\bibitem{AJ}
P.~Alexandersson and A.~Jal, \emph{Rook matroids and log-concavity of
$P$-Eulerian polynomials},
\href{https://arxiv.org/abs/2410.00127v3}{arXiv:2410.00127v3} (2026).

\bibitem{FMP}
L.~Ferroni, A.~H.~Morales and G.~Panova, \emph{Skew shapes, Ehrhart
positivity and beyond},
\href{https://arxiv.org/abs/2503.16403v3}{arXiv:2503.16403v3} (2026);
Proceedings of the London Mathematical Society \textbf{133} (2026),
e70204, \href{https://doi.org/10.1112/plms.70204}{doi:10.1112/plms.70204}.

\bibitem{JM}
K.~Jochemko and K.~Menon, \emph{Ehrhart positivity for marked order
polytopes}, \href{https://arxiv.org/abs/2604.08394v2}{arXiv:2604.08394v2}
(2026).

\bibitem{Kahane}
Y.~Kahane, \emph{Combinatorial interpretation of the coefficients of the
order polynomial of fence posets},
\href{https://arxiv.org/abs/2607.11225v1}{arXiv:2607.11225v1} (2026).

\bibitem{LXZ}
F.~Liu, G.~Xin and Z.~Zhang, \emph{Order polytopes of dimension $\le13$
are Ehrhart positive},
\href{https://arxiv.org/abs/2412.07164v1}{arXiv:2412.07164v1} (2024);
Discrete Mathematics (2025),
\href{https://doi.org/10.1016/j.disc.2025.114833}{doi:10.1016/j.disc.2025.114833}.

\bibitem{Stanley}
R.~P.~Stanley, \emph{Enumerative Combinatorics, Volume~1}, second
edition, Cambridge University Press, 2012.
\end{thebibliography}
\end{document}
